\uuid{pjiP}
\titre{Fonction croissante/décroissante}
\auteur{bodin ; hanani ; mzari}
\organisation{exo7}
\licence{CC BY-NC-SA 4.0}
\datecreate{2020-01-20}
\niveau{L1}
\module{Analyse}
\chapitre{Continuité, limite et étude de fonctions réelles}
\sousChapitre{Autre}
\difficulte{1}
\theme{Notion de fonctions, fonction croissante, décroissante}
\exo7qcmid{lille-1.continuite.001}

\contenu{
\begin{question}
\qclassification{123.00}
\qtags{motcle=fonction croissante/décroissante}
Quels arguments sont valides pour justifier que la fonction $x \mapsto \sin(x)$ n'est pas une fonction croissante sur $\Rr$ ?

\begin{answers}
    \bad{$\sin(\pi) = \sin(0)$ et pourtant $\pi \neq 0$.}

    \bad{$\sin(\frac\pi2) > \sin(0)$ et pourtant $0 < \frac\pi2$.}

    \good{$\sin(\frac{3\pi}{4}) > \sin(\pi)$ et pourtant $\frac{3\pi}{4} < \pi$.}

    \bad{On a $|\sin x| \le |x|$ pour tout $x\in\Rr$.}
\end{answers}
\begin{explanations}
Une fonction $f$ est croissante si $x \le y$ implique $f(x) \le f(y)$.
Donc une fonction n'est pas croissante si on peut trouver $x \le y$ mais avec
$f(x) > f(y)$. Le seul argument valable est donc $\frac{3\pi}{4} < \pi$ avec $\sin(\frac{3\pi}{4}) > \sin(\pi)$.
\end{explanations}
\end{question}
}
